\medskip

%\begin{tabularx}{0.8\linewidth}{|X|}\hline
%\textbf{Pour chaque question, recopier sur la copie son numéro et la réponse
%correspondante.}\\
%\textbf{Pour cette partie, aucune justification n’est demandée.}\\
%\textbf{Pour les questions a choix multiple, une seule réponse est exacte.}\\ \hline
%\end{tabularx}
%\end{center}

\medskip

\textbf{Question 1}

%Quel est le tiers de 18 ?
On a $\dfrac13 \times 18 = \dfrac{18}{3} = 6$.

\medskip

\textbf{Question 2}

%Un film dure $240$~min. Quelle est sa durée en heures ?
1 h égale 60 minutes, donc $\dfrac{240}{60} = 4$~(h).
\medskip

\textbf{Question 3}

Rangées dans l'ordre croissant les notes sont :

6~;~8~;~12~;~15~;~19 : la médiane est la troisième note soit 12.

\medskip

\textbf{Question 4}

%\begin{center}
%\psset{unit=1cm}
%\begin{pspicture}(0,-1)(9.,1)
%\psgrid[gridlabels=0pt,subgriddiv=4,gridwidth=0.25pt,subgridwidth=0.25pt]
%\psline(9.2,0)
%\qdisk(2,0){2pt}\qdisk(6,0){2pt}\qdisk(9,0){2pt}
%\uput[d](2,0){0}\uput[d](6,0){1}\uput[u](9,0){E}
%\end{pspicture}
%\end{center}
%
%Sur cette droite graduée, l’abscisse du point E est
%\begin{center}
%\begin{tabularx}{\linewidth}{*{4}{X}}
%\textbf{A.~} $\dfrac54$ &\textbf{B.~} $\dfrac32$ &\textbf{C.~} $\dfrac74$ &\textbf{D.~} $\dfrac52$\\ 
%\end{tabularx}
%\end{center}
En partant de l'origine point d'abscisse $0$, la point d'abscisse 1 est au 4\up{e} carreau, donc le point E a pour abscisse $\dfrac74$.
\medskip

\textbf{Question 5}

Les angles en A et C sont complémentaires (la somme en degré des trois angles vaut 180 et l'angle droit en B mesure 90\degres), on a donc $35 + \widehat{\text{C}} = 90$, soit $\widehat{\text{C}} = 90 - 35 = 55(\degres)$.
\medskip

\textbf{Question 6}
Par définition : $\widehat{\text{B}} = \dfrac{\text{mesure du côté adjacent}}{\text{mesure de l'hypoténuse}} = \dfrac{\text{BA}}{\text{BC}}$.

\medskip

\textbf{Question 7}

%\begin{minipage}{0.45\linewidth}
%Sur la figure ci-contre, dans le triangle ADE les droites (DE) et (CB) sont parallèles.
%
%Déterminer la longueur AD.
%\end{minipage}\hfill
%\begin{minipage}{0.52\linewidth}
%\psset{unit=0.9cm}
%\begin{pspicture}(8,3.5)
%\pspolygon(0.2,0.2)(7.7,0.2)(6.2,3.2)%AED
%\psline(2.4,0.2)(2,1.1)%BC
%\uput[ul](0.2,0.2){A} \uput[ur](7.7,0.2){E} \uput[ur](6.2,3.2){D} 
%\uput[d](2.4,0.2){B}\uput[ul](2,1.1){C}
%\uput[ur](2.2,0.6){2 cm} \uput[ul](1.2,0.7){4 cm} \uput[ur](6.95,1.7){7 cm}
%\end{pspicture}
%\end{minipage}
A, B et E sont alignés dans cet ordre ;

A, C et D sont alignés dans cet ordre ; et

les droites (DE) et (CB) sont parallèles.

On a donc une configuration de Thalès : les mesures des côtés des triangles ABC et AED sont proportionnelles, d'où en particulier $\dfrac{\text{AC}}{\text{AD}} = \dfrac{\text{AB}}{\text{AE}}$, soit avec les longueurs connues $\dfrac{4}{\text{AD}} = \dfrac27$, d'où $2\text{AD} = 4 \times 7$ ; on a donc $\text{AD} = 2 \times 7 = 14$~(cm).

\textbf{Question 8}

%Dans un collège, 25\,\% des 300 élèves participent à une olympiade de mathématiques.
%
%Combien d’élèves ne participent pas a cette olympiade ?
25\,\% des 300 élèves représentent $\dfrac{25}{100} \times 300 = \dfrac{25 \times 300}{100} = 25 \times 3 = 75$ (élèves).

$300 - 75 = 225$, donc 225 élèves ne participent pas à cette olympiade
\medskip


%\begin{minipage}{0.6\linewidth}
\textbf{Question 9}

%Une élève souhaite réaliser un programme avec un logiciel de
%programmation pour dessiner un carré.
%
%Par quelles valeurs doit-on compléter les lignes 3 et 5 
%pour obtenir un carré ?
%\end{minipage}\hfill
%\begin{minipage}{0.34\linewidth}
%\begin{scratch}[scale=0.8,num blocks, num start=1,baseline=2]]
%\initmoreblocks{definir carré}
%\blockpen{stylo en position d'écriture}
%\blockrepeat{répéter \ovalnum{\phantom{aa}} fois}
%{
%\blockmove{avancer de \ovalvariable{50} pas}
%\blockmove{tourner \turnleft{} de \ovalnum{\phantom{aa}} degrés}
%}
%\end{scratch}
%\end{minipage}
Il faut répéter 4 fois et tourner de 90 degrés.

%%%%%% Partie avec calculatrice
\begin{center}
\textbf{Partie 2 — Raisonnement et résolution de problèmes — 14 points — 1 h 40}
\end{center}

\medskip

%\begin{tabular}{|p{14cm}|}\hline
%\textbf{Dans cette partie, toutes les réponses doivent être justifiées, sauf si une indication
%contraire est donnée.}\\
%\textbf{La clarté et la précision des raisonnements ainsi que la rédaction sont évaluées sur
%2 points.}\\
%\textbf{Pour chaque question, si le travail n’est pas terminé, laisser tout de même une trace de la
%recherche ; les essais et les démarches engagées, même non aboutis, seront pris en
%compte dans la notation.}\\ \hline
%\end{tabular}

\bigskip
